\subsection{A partition of the spectrum}

Denote by \(\overline{Z}\) the subset \(Z \cup \{-\infty, +\infty\}\) of \(\overline{R}\) . If u is a linear map whose kernel or cokernel is finite-dimensional, one calls the index of u the element \(\text{ind}(u)\) of \(\overline{Z}\) defined by

\[
\operatorname{ind} (u) = \dim \operatorname{Coker} (u) - \dim \operatorname{Ker} (u)
\]

(cf. no. 6 of III, p. 67).

DEFINITION 3. --- Let E be a complete normable space and u an endomorphism of E. For every \(n \in \overline{Z}\) , denote by \(\mathrm{Sp}_{n}(u)\) the set of complex numbers \(\lambda \in \mathrm{Sp}_{\mathrm{e}}(u)\) such that \(u - \lambda1_{E}\) is a strict morphism, whose kernel or cokernel is finite-dimensional, and whose index is \(n\). We denote \(\mathrm{Sp}_{\omega}(u)\) the complement in \(\mathrm{Sp}_{\mathrm{e}}(u)\) of the union of the subsets \(\mathrm{Sp}_{n}(u)\) for \(n \in \overline{Z}\) .

The sets \(\mathrm{Sp}_{\mathrm{s}}(u)\) , \(\mathrm{Sp}_{n}(u)\) for \(n \in \overline{Z}\) and \(\mathrm{Sp}_{\omega}(u)\) form a partition of the spectrum of u.

Every endomorphism of E whose cokernel is finite-dimensional is strict (III, p. 52, lemma 6). The Fredholm endomorphisms of E are the endomorphisms of E whose kernel and cokernel are finite-dimensional (III, p. 52, prop. 11). The set \(\mathbf{C} - \mathrm{Sp}_{\mathrm{e}}(u)\) consists of the \(\lambda \in \mathbf{C}\) such that \(u - \lambda 1_{\mathrm{E}}\) is a Riesz endomorphism of E, and such an endomorphism is a Fredholm endomorphism of E of index 0.

Consequently, for \(\lambda \in \mathbf{C}\) and for \(n\in \mathbf{Z} - \{0\}\) , we have:

\(\lambda \in \mathbf{C} - \mathrm{Sp}(u) \iff u - \lambda 1_{\mathrm{E}}\) is an automorphism;

\(\lambda \in \mathrm{Sp}_{\mathrm{s}}(u) \quad \Longleftrightarrow \quad u - \lambda 1_{\mathrm{E}} \text{ is a Riesz endomorphism, but is not an automorphism;}\)

\(\lambda \in \mathrm{Sp}_{0}(u) \quad \Longleftrightarrow \quad u - \lambda 1_{\mathrm{E}}\) is a Fredholm endomorphism of index 0 of E but is not a Riesz endomorphism of E;

\(\lambda \in \mathrm{Sp}_n(u) \quad \Longleftrightarrow \quad u - \lambda 1_{\mathrm{E}}\) is a Fredholm endomorphism of index \(n\) of E (\(n \neq 0\));

\[
\lambda \in \mathrm{Sp} _ {- \infty} (u) \quad \Longleftrightarrow \quad u - \lambda 1 _ {\mathrm{E}} \text{ is strict, its kernel is infinite-dimensional and its cokernel is finite-dimensional;}
\]

\(\lambda \in \mathrm{Sp}_{+\infty}(u) \quad \Longleftrightarrow \quad u - \lambda 1_{\mathrm{E}}\) is strict, its kernel is finite-dimensional and its cokernel is infinite-dimensional;

\(\lambda \in \mathrm{Sp}_{\omega}(u) \quad \Longleftrightarrow \quad \text{either } u - \lambda 1_{\mathrm{E}} \text{ is not strict, or its kernel and cokernel are infinite-dimensional.}\)

Remark. --- Let \(\pi\) denote the canonical homomorphism of \(\mathcal{L}(\mathrm{E})\) onto the Calkin algebra \(\mathcal{C}alk(\mathrm{E})\) of E (cf. no. 4 of III, p. 75). The spectrum of \(\pi(u)\) with respect to the algebra \(\mathcal{C}alk(\mathrm{E})\) is sometimes called the stable spectrum of \(u\) . Its elements are the complex numbers \(\lambda\) such that \(u - \lambda 1_{\mathrm{E}}\) is not a Fredholm endomorphism of E (cor. 1 of th. 3 of III, p. 73). It is therefore equal to \(\mathrm{Sp}_{\omega}(u) \cup \mathrm{Sp}_{-\infty}(u) \cup \mathrm{Sp}_{+\infty}(u)\) .

Let A be the set of endomorphisms of E commuting with u. The closure B of \(\pi(\mathrm{A})\) in \(\mathcal{C}alk(\mathrm{E})\) is a complete normable algebra and the spectrum of \(\pi(u)\) with respect to B is the essential spectrum \(\mathrm{Sp}_{\mathrm{e}}(u)\) of u (III, p. 77, prop. 3). By the cor. of prop. 6 of I, p. 28, applied to the subalgebra B of \(\mathcal{C}alk(\mathrm{E})\) , the set \(\mathrm{Sp}_{\mathrm{e}}(u)\) is the union of \(\mathrm{Sp}_{\omega}(u) \cup \mathrm{Sp}_{-\infty}(u) \cup \mathrm{Sp}_{+\infty}(u)\) and of certain bounded connected components of its complement.

THEOREM 1. --- Let E be a complete normable space and u an endomorphism of E.

\begin{enumerate}
\def\labelenumi{\alph{enumi})}
\item
  The set \(\mathrm{Sp}_{\omega}(u)\) is compact. It is nonempty if E is of infinite dimension;
\item
  Let \(n \in \overline{Z}\) . The set \(\mathrm{Sp}_{n}(u)\) is a union of a family of bounded connected components of \(\mathbf{C} - \mathrm{Sp}_{\omega}(u)\) . It is open in C, and its boundary in C is contained in \(\mathrm{Sp}_{\omega}(u)\) .
\end{enumerate}

The set \(\mathbf{C} - \mathrm{Sp}_{\omega}(u)\) consists of the complex numbers \(\lambda \in \mathbf{C}\) such that \(u - \lambda 1_{\mathrm{E}}\) is a strict morphism whose kernel or cokernel is finite-dimensional. By Propositions 11 of III, p. 67 and 13 of III, p. 70, it is open. The set \(\mathrm{Sp}_{\omega}(u)\) is therefore closed. Since it is bounded, it is compact.

Let us prove \(b\) ). Let \(n \in \overline{\mathbf{Z}}\) . The set \(\mathrm{Sp}_n(u)\) is contained in \(\mathbf{C} - \mathrm{Sp}_{\omega}(u)\) . Let U be a connected component of \(\mathbf{C} - \mathrm{Sp}_{\omega}(u)\) which intersects \(\mathrm{Sp}_n(u)\) . The map \(\lambda \mapsto \operatorname{ind}(u - \lambda 1_{\mathrm{E}})\) of \(\mathbf{C} - \mathrm{Sp}_{\omega}(u)\) in \(\overline{\mathbf{Z}}\) being locally constant (cor. 1 of prop. 12 of III, p. 68 and cor. 1 of prop. 13 of III, p. 70), the index of \(u - \lambda 1_{\mathrm{E}}\) is equal to \(n\) for all \(\lambda \in U\) . If \(n \neq 0\) , this implies that U is contained in \(\mathrm{Sp}_n(u)\) . If \(n = 0\) , note that the set U is a connected component of \(\mathbf{C} - (\mathrm{Sp}_{\omega}(u) \cup \mathrm{Sp}_{-\infty}(u) \cup \mathrm{Sp}_{+\infty}(u))\) . Since U meets \(\mathrm{Sp}_0(u)\) and therefore \(\mathrm{Sp}_{\mathrm{e}}(u)\) , it follows from Remark 2 that the set U is contained in \(\mathrm{Sp}_{\mathrm{e}}(u)\) , and hence in \(\mathrm{Sp}_0(u)\) . In all cases, we conclude that \(\mathrm{Sp}_n(u)\) is the union of the connected components of \(\mathbf{C} - \mathrm{Sp}_{\omega}(u)\) that meet \(\mathrm{Sp}_n(u)\) . These are necessarily bounded since the set \(\mathrm{Sp}(u)\) is bounded. Consequently, \(\mathrm{Sp}_n(u)\) is open in \(\mathbf{C}\) and its boundary is contained in \(\mathrm{Sp}_{\omega}(u)\) . This proves \(b\) ).

Suppose finally that the set \(\mathrm{Sp}_{\omega}(u)\) is empty. By b), each of the sets \(\mathrm{Sp}_{n}(u)\) , for \(n \in \overline{Z}\) , is then empty. We therefore have \(\mathrm{Sp}(u) = \mathrm{Sp}_{\mathrm{s}}(u)\) . The spectrum of u is consequently discrete and compact, hence finite, and since all its points have finite spectral multiplicity, the vector space E is finite-dimensional (III, p. 83, prop. 2). This completes the proof of a).

COROLLARY. --- a) Let \(\Omega\) be the unbounded connected component of \(\mathbf{C} - \mathrm{Sp}_{\omega}(u)\) . We have \(\Omega \cap \mathrm{Sp}(u) \subset \mathrm{Sp}_{\mathrm{s}}(u)\) ;

\begin{enumerate}
\def\labelenumi{\alph{enumi})}
\setcounter{enumi}{1}
\tightlist
\item
  Every point adherent to \(\mathrm{Sp}_{\mathrm{s}}(u)\) that does not belong to \(\mathrm{Sp}_{\mathrm{s}}(u)\) belongs to \(\mathrm{Sp}_{\omega}(u)\) .
\end{enumerate}

The assertion \(a\) ) is a direct consequence of assertion \(b\) ) of th. 1. Let \(\lambda\) be a point adherent to \(\mathrm{Sp}_{\mathrm{s}}(u)\) that does not belong to \(\mathrm{Sp}_{\mathrm{s}}(u)\) . It belongs to the spectrum of \(u\) , since the latter is closed. It does not belong to any of the sets \(\mathrm{Sp}_{n}(u)\) , for \(n \in \overline{\mathbf{Z}}\) , since these are open (loc. cit.) and disjoint from \(\mathrm{Sp}_{\mathrm{s}}(u)\) . We therefore have \(\lambda \in \mathrm{Sp}_{\omega}(u)\) , whence \(b\) ).

PROPOSITION 3. --- Let E and F be complete normable spaces, \(u\colon E\to F\) and \(v\colon F\to E\) continuous linear maps.

\begin{enumerate}
\def\labelenumi{\alph{enumi})}
\item
  The traces on \(\mathbf{C} - \{0\}\) of the sets \(\mathrm{Sp}(v \circ u)\) and \(\mathrm{Sp}(u \circ v)\) (resp. \(\mathrm{Sp}_{\mathrm{s}}(v \circ u)\) and \(\mathrm{Sp}_{\mathrm{s}}(u \circ v)\) , resp. \(\mathrm{Sp}_{n}(v \circ u)\) and \(\mathrm{Sp}_{n}(u \circ v)\) for \(n \in \overline{\mathbf{Z}}\) , resp. \(\mathrm{Sp}_{\omega}(v \circ u)\) and \(\mathrm{Sp}_{\omega}(u \circ v)\) ) are equal;
\item
  Let \(\lambda\) be an element of \(\mathrm{Sp}_{\mathrm{s}}(v \circ u)\) distinct from 0. The spectral multiplicities of \(\lambda\) for \(v \circ u\) and for \(u \circ v\) are equal.
\end{enumerate}

Let \(\mu\) be a nonzero complex number and let \(n \in \overline{\mathbf{Z}}\) . For \(\mu 1_{\mathrm{E}} - v \circ u\) to be an automorphism (resp. a Riesz endomorphism, resp. a strict morphism whose kernel or cokernel is finite-dimensional and whose index is \(n\) ), it is necessary and sufficient that \(\mu 1_{\mathrm{F}} - u \circ v\) be one (III, p. 49, prop. 10). The assertion \(a\) ) then follows from the definitions.

Let \(\lambda\) be a point of \(\mathrm{Sp}_{\mathrm{s}}(v \circ u)\) distinct from 0. We have

\[
\dim \operatorname{Ker} ((\lambda 1 _ {\mathrm{E}} - v \circ u) ^ {n}) = \dim \operatorname{Ker} ((\lambda 1 _ {\mathrm{F}} - u \circ v) ^ {n})
\]

for all \(n \geqslant 0\) (loc. cit.), hence the spectral multiplicities of \(\lambda\) for \(v \circ u\) and \(u \circ v\) are equal.

